%% ============================================================
%% Paper04_COVID_Corruption_Summary.tex
%% One-page public summary
%% The Effect of the COVID-19 Pandemic on Risk of Corruption
%% Silverio-Murillo, Prudencio & Balmori-de-la-Miyar
%% Public Organization Review (2024)
%%
%% Compile: cd Slides
%%   TEXINPUTS=../Preambles:$TEXINPUTS xelatex -interaction=nonstopmode Paper04_COVID_Corruption_Summary.tex
%% ============================================================
\documentclass[9pt]{article}

\usepackage{geometry}
\geometry{a4paper, left=1.8cm, right=1.8cm, top=1.3cm, bottom=1.3cm}

\usepackage{fontspec}
\setmainfont{TeX Gyre Heros}
\setsansfont{TeX Gyre Heros}

\usepackage{xcolor}
\definecolor{navyblue}{HTML}{012169}
\definecolor{steelblue}{HTML}{2E75B6}
\definecolor{gold}{HTML}{F2A900}
\definecolor{medgray}{HTML}{555555}
\definecolor{lightbg}{HTML}{E8EDF5}
\definecolor{positivegreen}{HTML}{15803d}
\definecolor{negativered}{HTML}{B91C1C}

\usepackage{booktabs}
\usepackage{microtype}
\usepackage{multicol}
\setlength{\columnsep}{1.4em}
\usepackage{enumitem}
\setlist[itemize]{leftmargin=1.2em, itemsep=2pt, topsep=3pt, parsep=0pt}
\usepackage{parskip}
\setlength{\parskip}{3.5pt}
\setlength{\parindent}{0pt}
\usepackage{titlesec}
\titleformat{\section}[block]
  {\small\bfseries\color{navyblue}}
  {}{0em}{}
  [\vspace{1pt}{\color{steelblue}\hrule height 0.5pt}\vspace{2pt}]
\titlespacing*{\section}{0pt}{7pt}{1pt}

\pagestyle{empty}

%% ============================================================
\begin{document}

%% ---- Title block -------------------------------------------
{\centering
  {\large\bfseries\color{navyblue}
    The Effect of the COVID-19 Pandemic on Risk of Corruption}\\[0.4em]
  {\small\color{steelblue}\itshape
    Evidence from Mexico's Federal Procurement System}\\[0.3em]
  {\small\color{medgray}
    Adan Silverio-Murillo \enspace Daniel Prudencio \enspace
    Jose Roberto Balmori-de-la-Miyar\\
    \textit{Public Organization Review} 24 (2024) 1293--1313}\par
}

\vspace{0.4em}
{\color{navyblue}\hrule height 0.8pt}
\vspace{0.5em}

%% ---- Two-column body ---------------------------------------
\begin{multicols}{2}

\section{The Question}

Governments worldwide responded to COVID-19 by rapidly expanding
procurement budgets and relaxing accountability rules to speed up
purchasing. High-profile cases of overpriced medical supplies quickly
drew public attention to healthcare contracts.

We asked: did the pandemic actually increase \textbf{corruption risk}
in Mexico's federal procurement, and if so, \textbf{where did it
materialize?}

\section{Why It Matters}

Procurement corruption diverts public resources toward suppliers
chosen for political connections rather than cost and quality. During
a health emergency, when oversight is relaxed and budgets expand,
the window of opportunity widens.

Mexico's federal government allocates contracts through three
mechanisms: open public auctions, invitation-only auctions, and direct
allocation. The latter two are legally permitted but discretionary,
opening the door to favoritism. Before the pandemic, \textbf{70\% of
contract value} was already assigned through discretionary channels.

\section{The Data}

We use \textbf{CompraNet}, Mexico's official public procurement
platform, as systematized by the Instituto Mexicano para la
Competitividad (IMCO). The sample covers 64 federal institutions over
36 months, yielding 2,304 institution-month observations.

\vspace{0.4em}
{\small\centering
\begin{tabular}{lc}
  \toprule
  \textbf{Feature} & \textbf{Value} \\
  \midrule
  Federal institutions     & 64 \\
  Public acquisitions      & 378,000 \\
  Institution-month obs.   & 2,304 \\
  Period                   & Jan 2018 to Dec 2020 \\
  \midrule
  Pre-COVID DCVB (Overall) & 0.705 \\
  Pre-COVID DCVB (Health)  & 0.667 \\
  Pre-COVID DCVB (Non-Health) & 0.715 \\
  \bottomrule
\end{tabular}\par}

\vspace{0.4em}
The outcome variable is the \textbf{DCVB ratio}: the share of total
contract value assigned through discretionary mechanisms. A higher
DCVB signals greater corruption risk.

\section{How We Know}

The pandemic was a universal shock that affected all institutions at
the same time. We exploit the \textbf{time dimension} using a
difference-in-differences design: the COVID-19 indicator is one for
March through December 2020 and zero otherwise. Institution, month,
and year fixed effects absorb all time-invariant differences and common
seasonal patterns. An event study traces the effect month by month and
validates the parallel trends assumption.

\columnbreak

\section{What We Found}

The pandemic increased discretionary contracting by
\textbf{\textcolor{positivegreen}{+17\%}} on average. But the effect
is driven \textbf{entirely} by non-healthcare institutions.

\vspace{0.3em}
{\small\centering
\begin{tabular}{lcc}
  \toprule
  \textbf{Group} & \textbf{Effect} & \textbf{\% change} \\
  \midrule
  All institutions & \textcolor{positivegreen}{$+0.120^{***}$} & \textcolor{positivegreen}{$+17\%$} \\
  Healthcare       & $-0.060$ & $-9\%$ \\
  Non-healthcare   & \textcolor{positivegreen}{$+0.166^{***}$} & \textcolor{positivegreen}{$+23\%$} \\
  \bottomrule
\end{tabular}\par}

\vspace{0.4em}
Healthcare institutions, despite being at the center of the political
debate, show \textbf{no significant change} in corruption risk.
Non-healthcare institutions, such as defense, education, and
infrastructure, show a \textbf{23\% increase}.

The effect was also \textbf{temporary}: it peaked in May 2020
(two months after the stay-at-home order), and returned to pre-pandemic
levels within six months.

\section{How Did It Happen?}

Two mechanisms help explain the non-healthcare increase.

\textbf{Supplier concentration.} Non-healthcare institutions awarded
contracts to a smaller set of providers during the pandemic, as
measured by the Herfindahl-Hirschman Index. Fewer competing suppliers
means more room for favoritism.

\textbf{Reduced transparency.} The share of contracts whose procedural
information was published \textit{after} the contract start date rose
for non-healthcare institutions, a signal of reduced oversight.

Both mechanisms are consistent with a window of opportunity created
by weaker accountability, not a permanent shift in behavior.

\section{The Bigger Picture}

\begin{itemize}
  \item COVID-19 increased discretionary contracting by 17\%, but only
        in non-healthcare institutions, not in healthcare.
  \item Public scrutiny may itself be a corruption deterrent:
        sectors that received \textit{less} media attention are where
        risk actually materialized.
  \item Audit efforts targeting pandemic-era procurement should focus
        on non-healthcare institutions, even though the political
        debate centered on healthcare.
  \item DCVB measures corruption \textit{risk}, not actual corruption.
        The paper identifies a pattern at the aggregate level, not
        individual wrongdoing.
\end{itemize}

\end{multicols}

\vspace{0.2em}
{\color{navyblue}\hrule height 0.4pt}
\vspace{0.2em}
{\small\color{medgray}\centering
  Silverio-Murillo, Prudencio \& Balmori-de-la-Miyar (2024).
  ``The Effect of the COVID-19 Pandemic on Risk of Corruption in Public Procurement: Evidence from Mexico.''
  \textit{Public Organization Review} 24, 1293--1313.\par
}

\end{document}
